\documentclass[10pt,a4paper]{amsart}
\usepackage[margin=19mm]{geometry}
\usepackage{amsmath,amssymb,mathtools}
\usepackage[scr=rsfs,cal=euler]{mathalfa}
\usepackage{xfrac}
\usepackage[unicode,hidelinks,bookmarks=false]{hyperref}
\newcommand{\ie}{i.e.\ }
\newcommand{\cf}{cf.\ }
\newcommand{\resp}{respectively\ }
\newcommand{\Subsection}{\subsection}
\providecommand{\nobreakdash}{\nobreak\hbox{-}}
\setlength{\emergencystretch}{2em}
\multlinegap0pt
\usepackage{bm}
\usepackage{bbm}
\usepackage{dsfont}
\usepackage{xspace}
\usepackage{cancel}
\usepackage{mathtools}
\usepackage{amsthm}
\makeatletter
\@ifundefined{lem}{\newtheorem{lem}{Lemma}}{}
\makeatother
\title[Large gaps between values of several binary quadratic forms]{Large gaps between values of several binary quadratic forms}
\author{B{\l}a\.zej \.Zmija}
\date{Source: arXiv:2509.15365v1 (CC BY 4.0)}
\begin{document}
\maketitle

\noindent\emph{Source and licence.} Large gaps between values of several binary quadratic forms. Version arXiv:2509.15365v1 (CC BY 4.0), distributed under the Creative Commons Attribution 4.0 International licence, as recorded in the arXiv licence metadata for this exact version and shown on the abstract page https://arxiv.org/abs/2509.15365v1. Full rights notice: licenses/arxiv-2509.15365-CC-BY-4.0.txt.\par\smallskip

\noindent\textit{Editorial scope.} Counted unit: the Lem whose source label is \texttt{None} in the article (source0.tex lines 433--435, source label \texttt{None}), delivered together with the article's own complete original proof of it (source0.tex lines 436--466, one \texttt{proof} environment of 31 lines). No mathematics is written, repaired or re-derived: statement and proof are the article's own text line for line. Required context retained uncounted: LemPrime (lines 389--391). Every definition, display and label of those lines is retained unchanged. Omitted from this package and not redistributed: 1--388, 392--432, 467--626. Nothing inside a retained excerpt is omitted or altered: every retained body is delivered exactly as the article's lines are written, so no omission receipt of any kind accompanies this package. Citations: the retained excerpts carry no citation command at all, so no bibliography entry of the article is transcribed and no reference is redistributed. Reference adaptation: none was needed and none was made; every reference inside the delivered unit resolves inside it, so no unresolved reference or citation is produced. Printed slips kept literal: no printed word, symbol, hypothesis or number of the counted statement or of its proof was altered. Layout and class: the article's own class is \texttt{amsart} with packages including epsfig, url, setspace, color, babel, inputenc, fontenc, bm, bbold, amsmath, amsthm, amssymb, amsfonts, mathrsfs; the wrapper is the lane's pinned \texttt{amsart} editorial wrapper at 10pt on A4, which restates the theorem-like environments the retained text needs and adds the article's own macro definitions that the retained text needs (none), with extra wrapper packages (bm, bbm, dsfont, xspace, cancel, mathtools). The delivered numbering is the unit's own. Assets: no retained span contains a figure environment, a graphics-inclusion command, a PDF-page inclusion, a table or a TikZ picture; no image file is copied into this package, the package declares no asset dependency and no image file is redistributed with it. Licence: version arXiv:2509.15365v1 of this article is distributed under the Creative Commons Attribution 4.0 International licence, as recorded in the arXiv licence metadata for this exact version and shown on the abstract page https://arxiv.org/abs/2509.15365v1. A full rights notice is delivered at licenses/arxiv-2509.15365-CC-BY-4.0.txt. Characters: every retained excerpt is pure ASCII, so the delivered engine, XeLaTeX with its default Latin Modern text font, reports no missing character.
\begin{lem}\label{LemPrime}
Let $Q$ be a quadratic form with discriminant $D$. If $p$ is a prime and $\alpha$ is an odd number such that $\left(\frac{D}{p}\right)=-1$ and $p^{\alpha}\| n$, then $n$ is not represented by $Q$.
\end{lem}

\begin{lem}
For some $1\leq n \leq \delta L$ we have $\mathcal{S}\cap I_{n}=\emptyset$.
\end{lem}

\begin{proof}
Suppose that $m\in\mathcal{S}\cap I_{n}$ for some $n$. Then there exists $j\in\{1,\ldots,k\}$ such that $m=j+y_{0}+nP_{\delta}$. Hence,
\begin{align*}
dm\equiv d(j+y_{0})\equiv dj+dy\equiv dj+r \pmod{P_{\delta}}.
\end{align*}
However,
\begin{align*}
-1=\left(\frac{D}{r}\right)=\left(\frac{D}{dj+r}\right)
\end{align*}
for every $D\in\mathcal{D}$. Hence, there exists a prime $p$ and an odd number $\gamma$ such that $\left(\frac{D}{p}\right)=-1$ and $p^{\gamma}\| dj+r$. Then $\gamma\leq\beta_{p}$, or otherwise $p^{\gamma}>L\geq dj+r$. Write $dj+r=p^{\gamma}l$. Hence, if $\gamma\geq 3$ then $p^{3}\leq L$, so both $p\leq L/l_{t}$ and $p\leq\delta L$ are true if $L$ is big enough. Therefore, we get $p^{\beta_{p}+1}\mid P_{\delta}$. Hence, $p^{\gamma}\| m$, so $m\not\in\mathcal{S}$ by Lemma \ref{LemPrime}, a contradiction.

It follows from the previous paragraph that $\gamma =1$. We can also assume that $p>\delta L$ (otherwise $p\| m$, a contradiction by Lemma \ref{LemPrime}). Then $p\mid m$, and since $m\in\mathcal{S}$ and $\left(\frac{D}{p}\right)=-1$, we have $p^{2}\mid m$. Moreover, there exists $i\in\{1,\ldots, t\}$ such that $L/l_{i+1}<p\leq L/l_{i}$, so $l\leq l_{i}$ and hence $\varpi (p)\in T_{i}$, so $p\mid P_{\delta}$.

Now observe that for every pair $(j,p)$, where $j\leq k$ and $p>\delta L$, there is at most one $n<\delta L$ such that $p^{2}\mid j+y_{0}+nP_{\delta}$. Indeed, if it is true for numbers $n_{1}$ and $n_{2}$, then $p^{2}\mid (n_{1}-n_{2})P_{\delta}$, so $p\mid n_{1}-n_{2}$ because $p\| P_{\delta}$. However, $|n_{1}-n_{2}|\leq \delta L<p$, so $n_{1}=n_{2}$.

Moreover, if $p\mid j_{1}+y_{0}$ and $p\mid j_{2}+y_{0}$, then $p\mid j_{1}-j_{2}$. Hence, there are at most $\frac{L}{p}+1<\frac{1}{\delta}<\frac{2}{\delta}$ numbers $j$ such that $p\mid j+y_{0}$.

Let us summarize. So far we have proved that if $\mathcal{S}\cap I_{n}\neq \emptyset$ then there exist $j\leq L$ and a prime $p>\delta L$ such that $p\mid P_{\delta}$ and $p^{2}\mid j+y_{0}+nP_{\delta}$, and for every such a $p$ there are at most $\frac{2}{\delta}$ choices of $j$. Therefore,
\begin{align*}
\#\{n\leq\delta L\ |\ \mathcal{S}\cap I_{n}\neq\emptyset\} &\leq \#\big\{ (j,p)\ \big|\ j\leq L,\ p>\delta L,\ p\mid P_{\delta},\ p\mid j+y_{0}\big\}<\frac{2}{\delta}\#\{ p>\delta L\ |\ p\textrm{ divides } P_{\delta}\ \}.
\end{align*}
Define $F:=\#\{ p>\delta L\ |\ p\textrm{ divides } P_{\delta}\ \}$. Then 
\begin{align*}
(\delta L)^{F}\leq P_{\delta}<L^{2\pi (L)}<L^{2(1+\varepsilon )\frac{L}{\log L}}
\end{align*}
if $L$ is large enough, where $\pi (L)$ denotes the number of primes not greater than $L$. Therefore,
\begin{align*}
\#\{n\leq \delta L\ |\ \mathcal{S}\cap I_{n}\neq\emptyset\}<\frac{4(1+\varepsilon )}{\delta}\cdot\frac{L}{\log (\delta L)}.
\end{align*}
If $L$ is large enough, the last quantity is smaller than $\delta L$. The lemma is proved.
\end{proof}

\end{document}
